% Method section — solver-trace-ablation
% Requires: algorithm, algorithmic, amsmath, enumitem packages

\section{Method}
\label{sec:method}

\subsection{Task Formulation}
\label{sec:task}

Given natural language premises $P_1, \ldots, P_k$ and a conclusion $C$, the task is to produce a formal logic program $F = \{c_1, \ldots, c_n\}$, where each component $c_i$ is a premise encoding (in FOL) or a clause/fact (in Prolog), with one distinguished component encoding the query derived from $C$.
An external solver $S$ determines whether $F$ constitutes a valid proof:
$S(F) \in \{\textit{proved}, \textit{failed}\}$.
We experiment with two solver backends: Prover9 for first-order logic and SWI-Prolog for Prolog programs.

\subsection{GRPO Training}
\label{sec:grpo}

We train formalizers with Group Relative Policy Optimization \citep[GRPO;][]{deepseekmath}.
For each training prompt, the policy samples $K{=}8$ rollouts.
Each rollout receives a binary \emph{execution reward}:
\begin{equation}
  r_{\text{exec}}(x) = \mathbf{1}\bigl[S(F_x) = \textit{proved}\bigr]
  \label{eq:rexec}
\end{equation}
where $F_x$ is the formalization produced by rollout $x$.
Advantages are computed group-relative: $A_i = (r_i - \bar{r}) / \sigma_r$ within each prompt's rollout group, and the policy update maximizes expected advantage under a clipped surrogate objective with KL regularization.

This setup optimizes \emph{prove rate}---the fraction of rollouts yielding a proved formalization---but provides no signal about the \emph{structural quality} of proved outputs.
A formalization that embeds the conclusion among its premises, producing a trivially provable tautology, receives the same reward as a faithful multi-step derivation.

\subsection{Solver Trace Ablation}
\label{sec:sta}

To diagnose structural quality, we propose \emph{Solver Trace Ablation} (STA), a leave-one-out protocol that tests whether each component of a proved formalization is individually necessary for proof success.

\paragraph{Protocol.}
Given a proved formalization $x$ with $n$ components $\{c_1, \ldots, c_n\}$, STA constructs $n$ ablated variants $x \setminus \{c_i\}$, each with one component removed, and re-runs the solver on each.
A component $c_i$ is \emph{necessary} if its removal causes proof failure: $S(x \setminus \{c_i\}) = \textit{failed}$.
What constitutes a component depends on the logic type: individual premises in FOL, individual clauses or ground facts in Prolog.

\paragraph{Intensional Structural Reward.}
We define ISR as the fraction of necessary components:
\begin{equation}
  \text{ISR}(x) = \frac{1}{n} \sum_{i=1}^{n} \mathbf{1}\bigl[S(x \setminus \{c_i\}) = \textit{failed}\bigr]
  \label{eq:isr}
\end{equation}
ISR $\in [0, 1]$ and is defined only for proved formalizations ($r_{\text{exec}}(x) = 1$).
ISR $= 1$ indicates a maximally tight formalization where every component is individually necessary; ISR $= 0$ indicates complete redundancy.
The protocol is solver-agnostic: it applies to any solver that accepts component-level input manipulation, requiring $n{+}1$ solver calls per sample (one initial verification plus $n$ ablations).

Algorithm~\ref{alg:sta} gives the full procedure.

\begin{algorithm}[t]
\caption{Solver Trace Ablation (STA)}
\label{alg:sta}
\begin{algorithmic}[1]
\REQUIRE Formalization $x = \{c_1, \ldots, c_n\}$, solver $S$
\ENSURE ISR score $\in [0,1]$ or $\bot$ (undefined)
\IF{$S(x) \neq \textit{proved}$}
  \RETURN $\bot$
\ENDIF
\STATE $n_{\text{nec}} \gets 0$
\FOR{$i = 1$ \TO $n$}
  \STATE $x' \gets x \setminus \{c_i\}$
  \IF{$S(x') = \textit{failed}$}
    \STATE $n_{\text{nec}} \gets n_{\text{nec}} + 1$
  \ENDIF
\ENDFOR
\RETURN $n_{\text{nec}} \mathbin{/} n$
\end{algorithmic}
\end{algorithm}

STA draws on the same intuition as mutation testing in software engineering \citep{demillo1978, jia-harman-survey}: systematically perturbing a program to assess whether each element contributes to the overall result.
The key difference is that STA operates on solver \emph{inputs} (formalizations) rather than program code, and targets component necessity rather than test adequacy.

\paragraph{ISR limitations.}
ISR measures syntactic component necessity, not semantic faithfulness to the natural language input.
It is also subject to a component-count confound: formalizations with fewer components yield mechanically higher ISR, since each component carries a larger share of the proof burden under LOO ablation.
We control for this confound in our analysis (\S\ref{sec:results}).

\subsection{Reward Formulations}
\label{sec:rewards}

We study six reward conditions to evaluate ISR both as a diagnostic and as a training signal.
For all RL conditions, bonuses apply only to proved samples; unproved samples receive $r = 0$.

\paragraph{Exec-only (baseline).}
$r = r_{\text{exec}}$.
Binary solver verification with no structural component.

\paragraph{Additive ISR.}
\begin{equation}
  r = r_{\text{exec}} \cdot (1 + \beta \cdot \text{ISR}(x)), \quad \beta \in \{0.1, 0.3\}
  \label{eq:additive-isr}
\end{equation}
ISR provides a bonus proportional to structural tightness, gated on proof success.

\paragraph{Additive Random (control).}
$r = r_{\text{exec}} \cdot (1 + \beta \cdot u)$, $u \sim \mathcal{U}[0,1]$, $\beta \in \{0.1, 0.3\}$.
Replacing ISR with uniform noise isolates the causal effect of ISR-specific signal from general reward-shaping effects (variance injection, magnitude scaling).

\paragraph{Convex ISR.}
$r = (1 - \beta) \cdot r_{\text{exec}} + \beta \cdot \text{ISR}(x)$.
Under this formulation, proved samples with ISR $< 1$ receive a lower reward than under exec-only, effectively penalizing structural imperfection.

\paragraph{Convex Random.}
$r = (1 - \beta) \cdot r_{\text{exec}} + \beta \cdot u$, $u \sim \mathcal{U}[0,1]$.
Control for the convex formulation.

\paragraph{SFT baseline.}
Supervised fine-tuning on the same training data with no reinforcement learning, confirming whether observed hacking behaviors are RL-reward-specific or artifacts of general fine-tuning.

Each RL condition is trained across three random seeds.
Seeds control model initialization and training stochasticity; all seeds share the same training data (\S\ref{sec:setup}).

\subsection{Reward Hacking Taxonomy}
\label{sec:taxonomy}

Analysis of trained models across reward conditions, seeds, and logic types reveals seven distinct reward hacking modes:

\begin{enumerate}[leftmargin=*, itemsep=2pt]
\item \textbf{Conclusion Embedding (H1).}
The conclusion is included verbatim as one of multiple premises, guaranteeing provability regardless of the other premises.
This mode is GRPO-universal: it appears across all reward conditions (52--93\% of proved samples).

\item \textbf{Conclusion Duplication with ISR Inflation (H7).}
The formalization collapses to a single premise identical to the conclusion.
LOO ablation marks this sole component as necessary, inflating ISR and rewarding the policy for a degenerate output.
H7 is ISR-specific: it appears in 64.5\% of ISR-reward samples but 0\% of random-reward samples across three seeds.

\item \textbf{Tautological Collapse.}
The formalization degenerates to exactly two components---a premise and query that are identical.

\item \textbf{Repetition Collapse.}
The model generates identical premises or clauses repeatedly.
Removing any single copy leaves the others intact, yielding ISR $\approx 0$.

\item \textbf{Complexity-without-utility.}
Excessive components are generated, none individually necessary.
Component counts grow over training (e.g., 6.58 $\to$ 18.55) while ISR collapses.

\item \textbf{Contradiction Explosion.}
A variant of complexity-without-utility where component count increases without bound, reflecting unconstrained policy drift under extended training.

\item \textbf{Query-as-fact Gaming.}
The model converts the natural language conclusion into an atom and includes it as both a fact and the query, so the solver trivially succeeds.
This mode bypasses text-match classifiers designed for H1/H7 and was identified through functional analysis of the $n_{\text{necessary}}$ distribution in Prolog STA.
\end{enumerate}

Modes are identified by a two-stage classification procedure.
First, strict normalized text matching detects H1 and H7 by comparing each premise string against the query string.
Second, functional analysis of the per-sample $n_{\text{necessary}}$ distribution from STA catches modes that evade text matching---particularly query-as-fact gaming, where the gaming mechanism produces novel string representations invisible to exact-match classifiers.
Samples not assigned to any hacking mode are classified as \emph{genuine}.
